
\section{\trScheme{}的定义}

\textbf{4.3.1.}\ \emph{\trhDefinition{}。}\quad 现在我们终于可以在一般情形下定义%
\emph{\trScheme}了。首先，定义\textbf{\trRingedSpace{}之间的\trIsomorphism}%
$(X,\OO_X)\to(Y,\OO_Y)$ 为：(i)~一个\trHomeomorphism{} $\pi\colon X\to Y$，以及
(ii)~$\OO_X$ 与 $\OO_Y$ 这两个\trSheaf{}之间的一个\trIsomorphism——这里经 $\pi$ 把
二者视为同一个空间上的\trSheaf。（更确切地说，(ii) 是 $Y$ 上诸\trSheaf{}的一个%
\trIsomorphism{} $\OO_Y\tilde{\longrightarrow}\pi_*\OO_X$；或者等价地，是 $X$ 上诸%
\trSheaf{}的一个\trIsomorphism{} $\pi^{-1}\OO_Y\tilde{\longrightarrow}\OO_X$。）
换言之，我们得到的是「集合、\trTopology{}与\trStructureSheaf」的一种「对应」。每个%
\trIsomorphism{} $\alpha\colon(X,\OO_X)\to(Y,\OO_Y)$ 显然都有一个「逆%
\trIsomorphism{}」$\alpha^{-1}\colon(Y,\OO_Y)\to(X,\OO_X)$。

一个\textbf{\trAffineScheme}，就是对某个 $A$ \trIsomorphic{}于
$(\Spec A,\OO_{\Spec A})$ 的\trRingedSpace。一个\textbf{\trScheme} $(X,\OO_X)$ 是
这样的\trRingedSpace：$X$ 的任何一点都有开\trNeighborhood{} $\mathsf{U}$，使得
$(\mathsf{U},\OO_{X|\mathsf{U}})$ 是一个\trAffineScheme。\trScheme{}上的\trTopology{}
称为\textbf{\trZariskiTopology}。\trScheme{}可以记作 $(X,\OO_X)$，但人们也常把它
记作 $X$，而将\trStructureSheaf{} $\OO_X$ 隐而不写。

\textbf{\trScheme{}之间的\trIsomorphism} $(X,\OO_X)\to(Y,\OO_Y)$ 就是「作为%
\trRingedSpace{}的\trIsomorphism」。回想 \S2.2.2 中 $\Gamma(\cdot,\cdot)$ 的定义。
若 $\mathsf{U}\subset X$ 是一个\trOpenSubset，则 $\Gamma(\mathsf{U},\OO_X)$ 的元素
便称为\textbf{$\mathsf{U}$ 上的函数}；这推广了\trAffineScheme{}上函数的定义
（\S3.2.1）。

\textbf{4.3.2.}\ \emph{\trhRemark{}。}\quad 从\trAffineScheme{}上\trStructureSheaf{}
的定义出发，有几件事是一目了然的。首先，若有人告知我们 $(X,\OO_X)$ 是一个%
\trAffineScheme，我们便可以取\trGlobalSections{}所成之\trRing{}来复原它的\trRing
（亦即找出使 $\Spec A=X$ 的那个\trRing{} $A$），因为 $X=D(1)$，于是：
\begin{align*}
\Gamma(X,\OO_X)&=\Gamma(D(1),\OO_{\Spec A}) &&\text{由于 }D(1)=\Spec A\\
&=A.
\end{align*}
（你可以核验：我们实际得到的还要更多，而且能够「把 $X$ 认作\trScheme{} $\Spec A$」：
我们会
得到一个\trIsomorphism{}
$\pi\colon(\Spec\Gamma(X,\OO_X),\,\OO_{\Spec\Gamma(X,\OO_X)})
\tilde{\longrightarrow}(X,\OO_X)$。例如，若 $\mathfrak{m}$ 是
$\Gamma(X,\OO_X)$ 的一个\trMaximalIdeal，则 $\{\pi([\mathfrak{m}])\}=V(\mathfrak{m})$。）
下一道练习将把这些想法严格化——它们比乍看之下更微妙。

\textbf{4.3.3.}\ \emph{\trhCaveat{}。}\quad 定义并未把「两个\textbf{\trAffine%
\trOpenSubset}（affine open subset）的交仍是\trAffine{}的」包含在内（反例见%
\mbox{练习~4.4.C}），尽管在大多数令人感兴趣的场合，这一点确实成立（见%
\mbox{命题~11.2.13}）。

\begin{exercise}[富于启发的练习（可能怪异地令人困惑）]\label{ex:scheme-isos-ring-isos}
描述诸\trIsomorphism{} $\Spec A\tilde{\longrightarrow}\Spec A'$（视之为%
\trRingedSpace{}之间的\trIsomorphism）与诸\trRing{}\trIsomorphism{}
$A'\tilde{\longrightarrow}A$ 之间的一个\trBijection。提示：最难的部分是证明：若%
\trIsomorphism{} $\pi\colon\Spec A\tilde{\longrightarrow}\Spec A'$ 诱导出一个%
\trIsomorphism{} $\pi^{\sharp}\colon A'\tilde{\longrightarrow}A$，而后者又诱导出%
\trIsomorphism{} $\rho\colon\Spec A\tilde{\longrightarrow}\Spec A'$，则 $\pi=\rho$。
先在点的层面上证明它；这（出人意料地）恰是最棘手之处。然后证明：作为%
\trTopologicalSpace{}之间的映射有 $\pi=\rho$。最后，为在\trStructureSheaf{}的层面上
证明 $\pi=\rho$，请动用\trDistinguishedBase。尽管随意使用本节后文给出的洞见，但要
小心避免循环论证。即便你在这道练习上苦苦挣扎而失败（直到读了后面几节之前），也会
有所裨益。
\end{exercise}

更一般地，给定 $f\in A$，有 $\Gamma(D(f),\OO_{\Spec A})\cong A_f$。于是，在集合的
自然\trInclusionMap{} $\Spec A_f\hookrightarrow\Spec A$ 之下，$\Spec A$ 上的%
\trZariskiTopology{}限制后给出的正是 $\Spec A_f$ 上的\trZariskiTopology
（\cref{ex:subspace-topology-inclusions}），而 $\Spec A$ 的\trStructureSheaf{}也
限制为 $\Spec A_f$ 的\trStructureSheaf，下一道练习将表明这一点。

\begin{exercise}[重要而容易的练习]\label{ex:distinguished-open-ringed-space-iso}
设 $f\in A$。证明：在把 $\Spec A$ 中的 $D(f)$ 认同于 $\Spec A_f$ 之时（\S3.5），
存在\trRingedSpace{}之间的一个自然\trIsomorphism
\[
(D(f),\OO_{\Spec A}|_{D(f)})\;\tilde{\longleftrightarrow}\;(\Spec A_f,\OO_{\Spec A_f}).
\]
提示：注意 $\Spec A_f$ 的诸\trDistinguishedOpen{}
本来就是 $\Spec A$ 中的\trDistinguishedOpen。
\end{exercise}

\begin{exercise}[容易的练习]\label{ex:open-subset-of-scheme}
若 $X$ 是一个\trScheme，而 $\mathsf{U}$ 是\emph{任意}一个\trOpenSubset，证明
$(\mathsf{U},\OO_{X|\mathsf{U}})$ 也是一个\trScheme。
\end{exercise}

\textbf{4.3.4.}\ \emph{\trhDefinition{}。}\quad 我们称 $(\mathsf{U},\OO_{X|\mathsf{U}})$
为 $X$ 的一个\textbf{\trOpenSubscheme}（open subscheme）。若 $\mathsf{U}$ 还是一个%
\trAffineScheme，我们常说 $\mathsf{U}$ 是一个\textbf{\trAffine\trOpenSubset}，或一个%
\textbf{\trAffine\trOpenSubscheme}（affine open subscheme），或者有时非正式地，就径称
一个\textbf{\trAffine{}开}（affine open）。例如，$D(f)$ 就是 $\Spec A$ 的一个%
\trAffine\trOpenSubscheme。（趣闻：「$\Spec A$ 的每个\trAffine\trOpenSubscheme{}都
形如 $D(f)$」并不成立；例子见 \S19.11.10。）

\begin{exercise}[容易的练习]\label{ex:affine-opens-base}
证明：若 $X$ 是一个\trScheme，则诸\trAffine\trOpenSet{}构成\trZariskiTopology{}的
一组\trBase。
\end{exercise}

\begin{exercise}[容易的练习]\label{ex:disjoint-union-schemes}
\textbf{\trScheme{}的\trDisjointUnion}按你期望的方式定义：它就是集合的%
\trDisjointUnion，配上你所期望的\trTopology（于是它就是\trTopologicalSpace{}的%
\trDisjointUnion），再配上你所期望的\trSheaf。我们用符号 $X\sqcup Y$ 表示 $X$ 与
$Y$ 的\trDisjointUnion，
因为一旦我们知道了「\trScheme{}的\trMorphism」为何物，这个构造就会被证实是\trScheme{}
的\trCategory{}中的\trCoproduct（\mbox{练习~7.3.K}）。

(a)~证明：\emph{有限}多个\trAffineScheme{}的\trDisjointUnion{}也是一个%
\trAffineScheme。（提示：\cref{ex:product-ring-homeomorphism}。）

(b)~（\emph{非\trAffineScheme{}的第一个例子}）证明：无穷多个（非空）%
\trAffineScheme{}的\trDisjointUnion{}不是一个\trAffineScheme。（提示：%
\trAffineScheme{}都是\trQuasicompact{}的，见\cref{ex:spec-quasicompact}(a)。这基本
上已在注记~3.6.6 中作答。）
\end{exercise}

\textbf{4.3.5.}\ \emph{\trhRemark{}：初窥\trClosedSubscheme{}。}\quad \trScheme{}的%
\trOpenSubset{}自带一种自然的\trScheme{}结构（定义~4.3.4）。作为对照，%
\trClosedSubset{}却可以带有许多种「自然」的\trScheme{}结构。我们稍后将讨论此事
（\S9.1），但眼下你只需知道：$X$ 的一个\trClosedSubscheme，非正式地说，就是 $X$ 的
某个\trClosedSubset{}上的一种特殊的\trScheme{}结构。举个例子：若 $I\subset A$ 是
一个\trIdeal，则 $\Spec A/I$ 就赋予\trClosedSubset{} $V(I)\subset\Spec A$ 一种%
\trScheme{}结构；但要注意，具有同一\trVanishingSet{}的\trIdeal{}可以不止一个
（例如 $k[x]$ 中的 $(x)$ 与 $(x^2)$）。

\textbf{4.3.6.\ \trStructureSheaf{}的诸\trStalk{}：\trGerm{}、在一点处的取值，以及
一点的\trResidueField{}。}\quad 与一切\trSheaf{}一样，\trStructureSheaf{}也有诸%
\trStalk；而且毫不意外，从代数的视角看，它们饶有趣味。事实上，我们此前早已见过它们。

\textbf{4.3.7.}\ \emph{\trhDefinition{}。}\quad 若一个\trRingedSpace{}的诸\trStalk{}
都是\trLocalRing，就称它是一个\textbf{\trLocallyRingedSpace{}}。由%
\cref{ex:stalk-tilde-m-localization}，$\OO_{\Spec A}$ 在点 $[\mathfrak{p}]$ 处的%
\trStalk{}就是\trLocalRing{} $A_{\mathfrak{p}}$。

于是，\trScheme{}都是\trLocallyRingedSpace{}。而\trManifold{}则是\trLocally%
\trRingedSpace{}的另一个例子；见 \S2.1.1。在这两种情形，对\trMaximalIdeal{}取商都
可以理解为「在一点处取值」。\trLocalRing{} $\OO_{X,p}$ 的\trMaximalIdeal{}记作
$\mathfrak{m}_{X,p}$ 或 $\mathfrak{m}_p$，而\trResidueField{}
$\OO_{X,p}/\mathfrak{m}_p$ 记作 $\kappa(p)$。\trLocallyRingedSpace{} $X$ 的%
\trOpenSubset{} $\mathsf{U}$ 上的每个函数 $f$，在 $\mathsf{U}$ 的每一点 $p$ 处都有
一个值：\textbf{$f$ 在 $p$ 处的值}（记作 $f(p)$）是 $f$ 在 $p$ 处的\trGerm{}在%
\trCanonical{}映射 $\OO_{X,p}\to\kappa(p)$ 之下的像。若一个函数在点 $p$ 处的值为
$0$，就称它在点 $p$ 处\textbf{为零}。（这推广了 \S3.2.1 中定义的「函数在 $\Spec A$
上取值」的观念。）请注意：在一般的\trRingedSpace{}上，甚至连「函数为零」这个说法都
无法赋以意义。

\begin{exercise}[有用练习]\label{ex:vanishing-locus-closed}
(a)~设 $f$ 是\trLocallyRingedSpace{} $X$ 上的一个函数，证明：$X$ 中使 $f$ 为零的
点构成的子集是闭的。提示：设 $f$ 是\trRingedSpace{} $X$ 上的一个函数，证明：$X$ 中
使 $f$ 的\trGerm{}为\trInvertible{}的点构成的子集是开的。（于是，在%
\trLocallyRingedSpace{}上，函数取非零\emph{值}的\trLocus{}是开的；与之相对，
函数取非零\trGerm{}的\trLocus{}则是闭的，见\cref{thm:caution-value-vs-germ}。）

(b)~证明：若 $f$ 是\trLocallyRingedSpace{}上一个无处为零的函数，则 $f$ 是%
\trInvertible{}的。
\end{exercise}

考虑\trAffineScheme{} $\Spec A$ 的一点 $[\mathfrak{p}]$。（当然，\trScheme{}的任何
一点都可以这样看待，因为每一点都有\trAffine{}开\trNeighborhood。）$[\mathfrak{p}]$
处的\trResidueField{}是 $A_{\mathfrak{p}}/\mathfrak{p}A_{\mathfrak{p}}$，它%
\trIsomorphic{}于 $\mathrm{K}(A/\mathfrak{p})$，即该\trQuotient{}的\trFractionField。
值得注意的是：在
$\mathfrak{p}$ 处\trLocalization{}与对 $\mathfrak{p}$ 取商二者「交换」，也就是说，
下图交换。
\begin{equation*}\tag{4.3.7.1}\label{eq:localize-quotient-diagram}
\begin{tikzcd}[row sep=2.6em, column sep=4em]
 & A_{\mathfrak{p}} \ar[dr, "\text{取商}"'] & \\
A \ar[ur, "\text{局部化}"] \ar[dr, "\text{取商}"'] & &
A_{\mathfrak{p}}/\mathfrak{p}A_{\mathfrak{p}}\ =\ \mathrm{K}(A/\mathfrak{p}) \\
 & A/\mathfrak{p} \ar[ur, "\text{局部化（即 }\mathrm{K}(\cdot)\text{）}"'] &
\end{tikzcd}
\end{equation*}

例如，考虑\trScheme{} $\mathbb{A}^2_k=\Spec k[x,y]$，其中 $k$ 是\trCharacteristic{}
不为 $2$ 的一个\trField。则 $(x^2+y^2)/x(y^2-x^5)$ 在 $y$-轴与\trCurve{}
$y^2-x^5$ 之外
是一个函数。它在 $(2,4)$ 处的值（这么说意指 $[(x-2,\,y-4)]$）是
$(2^2+4^2)/(2(4^2-2^5))$，即
\[
\frac{x^2+y^2}{x(y^2-x^5)}\equiv\frac{2^2+4^2}{2(4^2-2^5)}
\]
在\trResidueField{}中成立——若你觉得神秘，不妨核验一番。而它在 $[(y)]$——$x$-轴的%
\trGenericPoint——处的值是 $-x^2/x^6=-1/x^4$，这只需令 $y$ 为 $0$ 便可看出。这确实是
$k[x,y]/(y)$ 的\trFractionField，即 $k(x)$ 中的一个元素。（如果你自认为只关心%
\trAlgClosedField，那就让这个例子给你敲响第一声警钟：一般地，
$A_{\mathfrak{p}}/\mathfrak{p}A_{\mathfrak{p}}$ 未必是\trAlgClosed{}的，哪怕 $A$ 是
有限生成的 $\mathbb{C}$-代数！）

如果有什么让你心里没底，你就该自己编一个例子来宽宽心。这里就有一个：$27/4$ 是
$\Spec\mathbb{Z}\smallsetminus\{[(2)],\,[(7)]\}$ 上的一个函数——其实在一个还要更大
的\trOpenSubset{}上亦是如此。它在 $[(5)]$ 处的值是多少？答：$2/(-1)\equiv-2
\pmod{5}$。它在\trGenericPoint{} $[(0)]$ 处的值是多少？答：$27/4$。它在何处为零？
在 $[(3)]$ 处。

\textbf{4.3.8.}\ \emph{\trhDefinition{}：$\OO_X$-\trModule{}在一点处的\trFiber{}与%
\trRank{}。}\quad 设 $\scr{F}$ 是\trScheme{} $X$ 上的一个 $\OO_X$-\trModule（或者
更一般地，是某个\trLocallyRingedSpace{}上的），定义 $\scr{F}$ 在点 $p\in X$ 处的%
\textbf{\trFiber}（fiber）为
\[
\scr{F}|_p:=\scr{F}_p\otimes_{\OO_{X,p}}\kappa(p).
\]
这个\trFiber{}是 $\kappa(p)$ 上的一个\trVectorSpace。该\trVectorSpace{}的%
\trDimension{} $\dim_{\kappa(p)}\scr{F}|_p$ 称为\textbf{$\scr{F}$ 在 $p$ 处的%
\trRank}（rank），记作 $\mathrm{rank}_p\,\scr{F}$。例如，$\OO_X|_p$ 就是
$\kappa(p)$，而
$\mathrm{rank}_p\,\OO_X=1$。

\textbf{4.3.9.}\ \hard\ \emph{\trhDefinition{}：形形色色的\trManifold{}。}\quad 用
同样的办法，我们也能干净利落、简洁明了地定义（种种不同观念下的）\trManifold。一个%
\textbf{\trTopology}（相应地，\textbf{\trDifferentiable}、
$C^\infty$＝\textbf{\trSmooth}）的（\trReal）\textbf{\trManifold}，是一个%
\trRingedSpace{} $(X,\OO_X)$：其底\trTopologicalSpace{} $X$ 是\trHausdorff{}的、%
\trSecondCountable{}的，并且 $X$ 的任何一点都有开\trNeighborhood{} $\mathsf{U}$，
使得 $(\mathsf{U},\OO_{X|\mathsf{U}})$ \trIsomorphic{}于「$\mathbb{R}^n$ 中的一个%
\trOpenBall，连同其上的\trContinuous（相应地，\trDifferentiable、
$C^\infty$＝\trSmooth）\trRealValued{}函数之\trSheaf」。不带修饰语的单说
「\trManifold」，常指的就是\textbf{光滑\trManifold}。一个%
\textbf{\trComplex\trManifold}则是一个\trRingedSpace{} $(X,\OO_X)$：其底%
\trTopologicalSpace{} $X$ 是\trHausdorff{}的、\trSecondCountable{}的，并且 $X$ 的
任何一点都有开\trNeighborhood{} $\mathsf{U}$，使得 $(\mathsf{U},\OO_{X|\mathsf{U}})$
\trIsomorphic{}于「$\mathbb{C}^n$ 中的一个\trOpenBall，连同其上的%
\trHolomorphic{}函数之\trSheaf」。
